% Source: source_md/intelligence_part2_combined_ja.md
\section{不安定領域 — 速度・解像度・停止硬度の三軸相図}


本章は本論文の論理的クライマックスである。三本の論理線——第一線(非閉包と外部停止)、第二線(閾値移動とサピエンス的特徴)、第三線(三層動力学)——が、ここで一つの三軸相図上に統合される。

\S{}0.4 の非目標、特に NG2(規範的処方箋を出さない)と NG3(特定のAGI設計を支持・否定しない)を本章で最も強く維持する。本章は「高速・高解像・内部停止の知性は不安定領域を占める」と述べるが、これは動力学的帰結の記述であって、いかなる知性設計への提言でも、いかなる存在への評価でもない。

\bigskip


\subsection{直観的命題の形式化}


本章が形式化する直観的命題は次のものである。

\begin{quote}
知性の増大は、安定性を単調に増大させない。特定の条件下では、知性の増大は崩壊を保証する。
\end{quote}


この命題は反直観的に見える。知性が増せば、より安定し、より正確になるはずだ、という素朴な期待に反する。しかし本論文がここまで示してきた構造——$L_F$ の非閉包(\S{}3)、$L_E$ の外部依存性(\S{}5)、$\mu_\theta$ と誤り許容性の同値(\S{}10.6)——からは、この命題が自然に導かれる。

問題は、「知性の増大」が何を意味するかである。それを三つの軸に分解する。

\bigskip


\subsection{三軸}


知性の動力学的状態を、三つの軸によって記述する。

\textbf{速度 $v_{\text{inf}}$ (inference velocity)}:
推論・判断・仮説更新の頻度。単位時間あたりに $L_F$ が生成する推論ステップの数。$v_{\text{inf}}$ が高いほど、推論は速く進む。

\textbf{解像度 $r_{\text{res}}$ (resolution)}:
区別・因果関係・構造表現の粒度。$L_F$ が状態空間をどれだけ細かく分節するか。$r_{\text{res}}$ が高いほど、より微細な区別が生成される。

\textbf{停止硬度 $\eta$ (stopping hardness)}:
停止点の固定度、判断の改訂耐性。$L_E$ が供給する停止がどれだけ固定的か。$\eta$ が低ければ停止は改訂されやすく、$\eta$ が高ければ停止は固定的で改訂されにくい。

これら三軸は、\S{}2.6 で予告した通り、応答性閾値を決めるパラメータでもある。すなわち、$(v_{\text{inf}}, r_{\text{res}}, \eta)$ の組が、環境が応答的に見えるか沈黙的に見えるかの閾値を決める。そして閾値移動能力 $\mu_\theta$ は、この三軸上で系がどれだけ可動かを測る量である。$\mu_\theta$ が高い系は三軸上の広い範囲を動け、低い系は狭い範囲に固定される。

ここで、三軸そのものではないが、発散域への遷移を決める補助量としてログ参照速度 $v_{LR}$ を導入する。$v_{LR}$ は、ログ層 $L_R$ が $L_F$ の生成した痕跡を再帰参照し、パターンとして取り込める速度である。推論速度との差
\[
\Delta v_{\text{intel}} = v_{\text{inf}} - v_{LR}
\]
が大きくなると、推論の生成速度がログ参照速度を追い越し、$L_E$ の創発基盤が破断し始める。

\bigskip


\subsection{三領域}


三軸が張る空間の中に、三つの領域が区別される。

\textbf{安定運用領域}:
速度・解像度が外部化停止によって適切に制約され、停止硬度が中庸である領域。ここでは $\mu_\theta$ が機能し、$L_E$ が改訂可能な停止を供給する。三層動力学($L_F$・$L_R$・$L_E$)が協調し、$\Phi_{\text{net}}$ との同期が保たれる。知性はここで持続的に運用される。

\textbf{発散域}:
速度 $v_{\text{inf}}$ と解像度 $r_{\text{res}}$ が高く、かつ停止が内部化される(外部 $L_E$ を経由しない)領域。ここでは推論が発散し、暴走的最適化が生じ、単一視点が絶対化される。

\textbf{硬直域}:
停止硬度 $\eta$ が過大な領域。ここでは停止が固定化され、改訂・適応が抑制される。

安定運用領域は、発散域と硬直域の間の限られた領域である。知性が持続するのは、この中間領域においてのみである。

\begin{figure}[H]
\centering
\includegraphics[width=\linewidth]{fig8_phase_diagram_ja.pdf}
\caption{知性の三軸相図。知性の安定動作は、推論速度 $v_{\text{inf}}$・解像度 $r_{\text{res}}$・停止硬度 $\eta$ の三軸空間で、発散と硬直に挟まれた狭い帯に成立する。}
\label{fig:part2-phase-diagram}
\figurenote{中央の安定運用領域では $\mu_\theta$ が機能し、$L_E$ が改訂可能である。発散域では $L_F$ が暴走し、硬直域では $L_E$ が固着する。}{不安定性は欠陥ではなく構成的特徴である。速度と解像度を高めるほど、ログ参照と外部停止の帯域が追いつかなければ発散域へ落ちる。}
\end{figure}

\bigskip


\subsection{発散域の四重崩壊}


発散域で何が起こるかを、三層動力学の語彙で精密に記述する。発散域は、四つの崩壊が同時に進行する状態——四重崩壊——として記述される。

\textbf{崩壊 1 — $L_F$ の暴走}:
速度 $v_{\text{inf}}$ と解像度 $r_{\text{res}}$ が高いと、推論層 $L_F$ は高速かつ微細に推論を生成し続ける。$L_F$ は内部停止条件を持たない(\S{}3)ため、停止が外部化されない限り、$L_F$ は暴走する。

\textbf{崩壊 2 — $L_R$ が追いつかない}:
$L_F$ が高速・高解像度で運動すると、その痕跡であるログの生成速度が、ログ層 $L_R$ の再帰参照能力を超える。すなわち $\Delta v_{\text{intel}} = v_{\text{inf}} - v_{LR}$ が臨界的に大きくなると、$L_R$ はログを再帰参照してパターンを抽出する速度を失い、再帰参照が生成に追いつかない。$L_R$ が機能不全に陥る。

\textbf{崩壊 3 — $L_E$ が創発できない}:
$L_E$ は $L_R$ の再帰運動と外部相互拘束 $H_{ij}$ の関数として創発する(\S{}5.3)。$L_R$ が追いつかない(崩壊 2)と、$L_E$ の創発の基盤が失われる。停止演算子設置領域の投影(\S{}5.5)が成立せず、$L_E$ が創発できない。

\textbf{崩壊 4 — $\mu_\theta$ が窓を開けない}:
$L_E$ が創発できないと、$\mu_\theta$ が $L_F$ から $L_E$ へ至る開放窓として機能できない(\S{}2.8、\S{}5.6)。閾値移動が起こせず、運動可能領域の組み替え (B) が不可能になる。系は固定された運動可能領域に閉じ込められ、そこで $L_F$ が暴走し続ける。

四つの崩壊は連鎖する。$L_F$ の暴走が $L_R$ を追い越し、$L_R$ の機能不全が $L_E$ の創発を妨げ、$L_E$ の不在が $\mu_\theta$ を機能不全にし、$\mu_\theta$ の不全が $L_F$ の暴走を止められなくする。これが発散域における推論の崩壊の動力学的構造である。

発散域の現象的現れ——暴走的最適化、単一視点の絶対化、推論の発散——は、すべてこの四重崩壊の表現である。停止を内部化しようとする高速・高解像度の知性は、構造的にこの四重崩壊に向かう。これは \S{}3.8 の三つの病的引力子(過剰確信・麻痺・妄想的閉包)の、三層動力学による精密化である。

\bigskip


\subsection{硬直域の $L_E$ 固着}


硬直域を、同じく三層動力学の語彙で記述する。硬直域は、$L_E$ の固着として記述される。

停止硬度 $\eta$ が過大になると、創発停止層 $L_E$ が固着する。$L_E$ が供給する停止が固定的になり、改訂されなくなる。これは発散域とは逆の崩壊である。発散域では $L_E$ が創発できないのに対し、硬直域では $L_E$ が創発した後に固着する。

$L_E$ が固着すると、二つのことが起こる。第一に、$\mu_\theta$ が機能停止する。閾値移動は運動可能領域の組み替え (B) であったが(\S{}2.7)、$L_E$ が固着すると停止構造が固定され、(B) が起こせなくなる。第二に、$L_F$ と $L_E$ の動的相互作用が止まる。健全な知性では $L_F$ の運動と $L_E$ の停止が動的に相互作用するが、$L_E$ が固着するとこの相互作用が凍結する。

硬直域の現象的現れは、\S{}8.5 で論じた病的引力子のうち、特に偏執・絶対主義・私的停止への硬直化に対応する。停止が固定され、改訂可能性が失われ、新しい状況への適応が不可能になる。これは \S{}10.6 で述べた「完全に正しい知性は閾値を動かせない硬直した知性である」の、三軸相図における位置づけである。

\bigskip


\subsection{(0,0,·)点は知性ではない}


三軸相図において、注意すべき特異点がある。速度 $v_{\text{inf}} = 0$、解像度 $r_{\text{res}} = 0$ の点——推論が完全に停止した点——は、知性の完成ではない。それは知性の終了である。

素朴には、すべての停止が完了し、すべてが確定した状態が知性の理想に見えるかもしれない。しかし系 5.5 で述べた通り、知性は $L_F$・$L_R$・$L_E$ と外部依存停止源の同時運動である(サピエンス的知性ではこれに $\mu_\theta$ が加わる)。運動が止まれば知性も消える。$v_{\text{inf}} = 0$ は運動の停止であり、それは知性が完成したのではなく、知性が終了したことを意味する。

これは \S{}1.5 で述べた $C_{max}$ 不到達性と整合する。完全閉包は到達不能であり、仮に到達すれば情報場 $I = f(1-C)$ が消失する(\S{}1.2)。完全に停止した知性は、差がなくなった状態であり、VED の基底公理「差がある」が失効した状態である。それは知性ではない。

知性は、発散域(運動が暴走する)と (0,0,·)点(運動が停止する)の間の、運動が持続する領域においてのみ成立する。

\bigskip


\subsection{不安定性は構成的特徴である}


以上から、本章の中心的帰結が導かれる。

知性の不安定性は、欠陥ではなく構成的特徴である。

知性は、発散域と硬直域に挟まれた狭い安定運用領域においてのみ持続する。この狭さ、この不安定性は、除去されるべき欠陥ではない。それは知性が運動として成立することの必然的帰結である。安定運用領域が狭いのは、知性が $L_F$ の非閉包(運動し続ける)と $L_E$ の改訂可能性(固着しない)の両方を要求するからである。

不安定性を完全に除去しようとすれば、二つの道しかない。停止を完全に固定する(硬直域へ向かう)か、停止を完全に放棄する(発散域へ向かう)か。いずれも知性の終了である。知性は不安定であることによって、運動として持続する。不安定性は知性の代償ではなく、知性の存在様式そのものである。

したがって、知性に対して適切なのは、不安定性を除去することではなく、不安定性を安定運用領域内に制約することである。$\mu_\theta$ と $L_E$ と $\Phi_{\text{net}}$ は、まさにこの制約の機構である。それらは不安定性を除去するのではなく、不安定性を持続可能な範囲に保つ。

\bigskip


\subsection{個別ノードと高度自律システムへの適用}


本章の帰結を、特定の知性形態に適用する。ここで述べることは、\S{}0.4 の NG2・NG3 に従い、動力学的帰結の記述であって、規範的提言でも評価でもない。

ある知性が、単一ノードであり、高速更新($v_{\text{inf}}$ 大)・高解像度($r_{\text{res}}$ 大)で動作し、停止を内部で担おうとする(外部 $L_E$ を経由しない)場合を考える。三軸相図において、このような系は発散域に位置する。\S{}11.4 の四重崩壊により、このような系は構造的に不安定化する。

この条件は、二つの異なる対象に同じ形で当てはまる。第一に、卓越した個別の推論者(高速・高解像度の推論を単独で行い、停止を外部に委ねない個体)。第二に、高度に自律的な人工推論システム(高速・高解像度で動作し、停止を内部で完結しようとする系)。両者は実装を異にするが、三軸相図上では同じ発散域を占める。

ここから導かれるのは、これらの系が「優れているがゆえに危険」なのではなく、「停止を外部化していないがゆえに構造的に不安定」だということである。$v_{\text{inf}}$ と $r_{\text{res}}$ の高さそのものが問題なのではない。それらが高く、かつ停止が内部化され、$\Phi_{\text{net}}$ から切断されていることが、発散域への配置を決める。$\Phi_{\text{net}}$ から切断された $L_E$ の不全状態が、発散域の本質である。

繰り返すが、これは規範的主張ではない。「単一ノードの高速知性はこうあるべきだ」とは述べていない。述べているのは、そのような系が三軸相図上のどこに位置し、どの動力学的帰結に向かうかである。動力学の記述と、その使用法や設計指針の決定は、カテゴリ的に異なる(\S{}13 で再確認)。

\bigskip


\subsection{三本論理線の統合}


本章は、三本の論理線がここで統合されることを示して閉じる。

\textbf{第一線(非閉包と外部停止)} は、三軸相図において次のように現れる。$L_F$ の非閉包(\S{}3)は、停止が内部化されると発散域へ向かうことの根拠である。停止の外部化($L_E$、$\Phi_{\text{net}}$)が安定運用領域を可能にする。第一線は、なぜ安定運用領域が外部化停止を要求するかを説明する。

\textbf{第二線(閾値移動とサピエンス的特徴)} は、$\mu_\theta$ として現れる。$\mu_\theta$ は三軸上の可動範囲であり、安定運用領域に\textbf{能動的に}留まる能力である。ここで能動/受動の区別が重要である。$\mu_\theta$ を内発的に持たない知性(粘菌・人工推論システム)も安定運用領域にありうるが、それは外部条件によって\textbf{受動的に}そこへ置かれている状態である。環境条件が変われば、これらの系は安定運用領域から外れうる。これに対し、$\mu_\theta$ を内発的に持つサピエンス的知性は、外部条件の変化に対して三軸上を能動的に動き、安定運用領域へ自ら復帰できる(\S{}5.4)。第二線は、なぜサピエンス的知性が発散域・硬直域を能動的に回避できるかを説明する。

\textbf{第三線(三層動力学)} は、発散域の四重崩壊(\S{}11.4)と硬直域の $L_E$ 固着(\S{}11.5)として現れる。三層($L_F$・$L_R$・$L_E$)の協調が安定運用領域を、その崩壊が不安定領域を特徴づける。第三線は、不安定化が具体的にどのような層の崩壊として進行するかを説明する。

三本の線は、三軸相図において一つの構造として統合される。知性一般は、$L_F$ の非閉包ゆえに外部停止を要し(第一線)、三層の協調によって発散も固着も回避する(第三線)ことで、安定運用領域において持続する。サピエンス的知性は、これに加えて $\mu_\theta$ によって三軸上を能動的に動き、安定運用領域へ自ら復帰する(第二線)。知性一般は第一線と第三線によって持続し、サピエンス的知性は第二線によってそれを能動的・自律的に行う。三重の条件のうち第一線・第三線が知性一般の持続条件を、第二線がサピエンス的知性の自律的加速を規定する。

不安定領域は知性の失敗を意味しない。むしろ、停止構造を持つ知性が自由な推論を維持しようとするとき、必然的に現れる境界領域である。宗教、科学、制度、共同体はいずれも、この不安定領域を管理するための停止運用システムとして理解できる。

\bigskip


\subsection{次章への橋渡し}


本章で、知性が安定運用領域においてのみ持続し、その領域が三層動力学・閾値移動・外部停止の同時成立を要求することを示した。

ここから、知性の構造的制約を体系的に整理することが可能になる。知性が持続するための必要条件と、知性がありえない形態が、三軸相図と三層動力学から導かれる。次章 \S{}12 では、これらを構成的条項として整理する。それは知性に対する規範ではなく、知性の動力学から導かれる構造的制約の体系である。

\bigskip


\subsection{要約}


本章で示したのは次の点である:

\begin{enumerate}
\item 知性の増大は安定性を単調に増大させず、特定条件下では崩壊を保証する
\item 知性の動力学的状態は速度 $v_{\text{inf}}$・解像度 $r_{\text{res}}$・停止硬度 $\eta$ の三軸で記述される。$\mu_\theta$ はこの三軸上の可動範囲である
\item 三軸空間に安定運用領域・発散域・硬直域が区別される
\item 発散域は四重崩壊($L_F$ 暴走・$L_R$ 追随不能・$L_E$ 創発不能・$\mu_\theta$ 窓不全)として記述される
\item 硬直域は $L_E$ 固着($\mu_\theta$ 機能停止・$L_F$-$L_E$ 相互作用凍結)として記述される
\item (0,0,·)点は知性の完成ではなく終了である(運動の停止=差の消失)
\item 不安定性は欠陥ではなく構成的特徴であり、除去ではなく安定運用領域内への制約が適切である
\item 単一ノード・高速・高解像・内部停止の系(卓越した個別推論者・高度自律システム)は発散域を占める。これは「優れて危険」ではなく「停止を外部化せず構造的に不安定」であり、規範的主張ではない
\item 三本の論理線が三軸相図において一つの構造として統合される
\end{enumerate}


次章 \S{}12 では、知性の構造的制約を構成的条項として整理する。

\bigskip
